\documentclass{uai2026} % for initial submission

%% Choose your variant of English; be consistent
\usepackage[american]{babel}

%% Packages
\usepackage{natbib}
    \bibliographystyle{plainnat}
    \renewcommand{\bibsection}{\subsubsection*{References}}
\usepackage{mathtools}
\usepackage{booktabs}
\usepackage{multirow}
\usepackage{graphicx}
\usepackage{algorithm}
\usepackage{algorithmic}
\usepackage{tikz}

%% Macros
\newcommand{\ie}{i.e.}
\newcommand{\eg}{e.g.}

\title{Conformal Prediction Under Distribution Shift: A COVID-19 Natural Experiment}

% Author info hidden for anonymous submission (automatically by class)
\author[1]{Chorok Lee}
\affil[1]{%
    Korea Advanced Institute of Science and Technology (KAIST)\\
    Daejeon, South Korea
}

\begin{document}
\maketitle

\begin{abstract}
Conformal prediction guarantees degrade under distribution shift.
We study this using COVID-19 as a natural experiment across 8 supply chain tasks.
Despite identical severe feature turnover (Jaccard approximately 0), coverage drops vary
from 0\% to 86.7\%, spanning two orders of magnitude. Using SHapley Additive exPlanations
(SHAP) analysis, we find catastrophic failures correlate with single-feature
dependence ($\rho = 0.714$, $p = 0.047$). Catastrophic tasks concentrate importance
in one feature (4.5$\times$ increase), while robust tasks redistribute across many
(10--20$\times$). Quarterly retraining restores catastrophic task coverage from 22\%
to 41\% (+19 pp, $p = 0.04$), but provides no benefit for robust tasks (99.8\%
coverage). Exploratory analysis of 4 additional tasks with moderate feature stability
(Jaccard 0.13--0.86) reveals feature stability, not concentration, determines
robustness, suggesting concentration effects apply specifically to severe shifts.
We provide a decision framework: monitor SHAP concentration before deployment;
retrain quarterly if vulnerable ($>$40\% concentration); skip retraining if robust.
\end{abstract}

\section{Introduction}\label{sec:intro}

Conformal prediction provides distribution-free coverage guarantees under the assumption of exchangeability~\citep{vovk2005algorithmic}. However, real-world deployments face distribution shifts that violate this assumption. A critical open question is: \textbf{How do conformal guarantees degrade under distribution shift, and what factors determine the severity of this degradation?}

We leverage a supply chain benchmark dataset with temporal splits aligned to COVID-19 onset (February 2020) and peak (July 2020) to conduct a controlled study of conformal prediction under distribution shift. Our key contributions:

\begin{enumerate}
    \item \textbf{Quantification}: We measure coverage degradation across 8 supply chain tasks experiencing severe feature turnover (Jaccard $\approx$0), finding drops ranging from 0\% to 86.7\%---two orders of magnitude variation under identical temporal shift.
    \item \textbf{Predictive Signal}: Using SHAP analysis, we identify single-feature dependence as a predictive signal for catastrophic failure in severe-shift scenarios (Spearman $\rho = 0.714$, $p = 0.047$, $n=8$). Exploratory analysis of 4 additional tasks with moderate feature stability (Jaccard 0.13--0.86) reveals a different mechanism.
    \item \textbf{Practical Solution}: Quarterly retraining significantly restores catastrophic task coverage by 19 percentage points (Wilcoxon $p=0.04$ vs no retraining).
    \item \textbf{Cost-Benefit}: Robust tasks (99.8\% coverage) gain no benefit from retraining, enabling practitioners to avoid unnecessary computational cost.
    \item \textbf{Decision Framework}: Actionable guidance based on SHAP importance concentration: monitor before deployment, retrain quarterly if vulnerable ($>$40\% concentration), skip retraining if robust.
    \item \textbf{Negative Result}: Adaptive Conformal Inference does not help under severe distribution shift.
\end{enumerate}

\section{Related Work}\label{sec:related}

\textbf{Conformal Prediction.} \citet{vovk2005algorithmic} introduced conformal prediction with exchangeability guarantees. Recent work extends to classification~\citep{romano2020classification} and regression~\citep{romano2019conformalized}. \citet{angelopoulos2021gentle} provide a comprehensive tutorial.

\textbf{Distribution Shift.} \citet{tibshirani2019conformal} study conformal prediction under covariate shift with known propensity scores. \citet{podkopaev2021distribution} develop distribution-free prediction sets under label shift. \citet{barber2023conformal} provide a comprehensive treatment beyond exchangeability. We focus on temporal shift with feature staleness---a distinct failure mode where feature distributions have zero overlap between train and test.

\textbf{Distribution Shift Benchmarks.} WILDS~\citep{koh2021wilds} provides standardized evaluation for in-the-wild shifts, while Shifts~\citep{malinin2021shifts} focuses on safety-critical applications. \citet{gulrajani2020search} show that many ``robust'' methods fail under temporal shift. We focus specifically on how conformal prediction guarantees degrade.

\textbf{Adaptive Methods.} \citet{gibbs2021adaptive} propose Adaptive Conformal Inference (ACI) for non-stationary settings, with extensions by \citet{zaffran2022adaptive} for time series. We show that these methods fail when feature overlap is $\sim$0\%.

\textbf{Interpretability.} Our finding that SHAP importance concentration predicts conformal failure connects to work on feature attribution~\citep{lundberg2017unified} and model debugging~\citep{adebayo2018sanity}. We show SHAP can serve as a pre-deployment diagnostic.

\section{Methodology}\label{sec:method}

\subsection{Dataset and Temporal Splits}

We use a supply chain dataset with temporal splits:
\begin{itemize}
    \item \textbf{Train}: Before February 2020 (pre-COVID)
    \item \textbf{Validation}: February--July 2020 (COVID onset)
    \item \textbf{Test}: After July 2020 (COVID peak)
\end{itemize}
This provides a natural experiment where COVID-19 serves as a documented distribution shift event.

\subsection{Conformal Prediction Setup}

We use conformal prediction via the Adaptive Prediction Sets (APS) algorithm~\citep{romano2020classification} with $\alpha = 0.1$ (90\% target coverage):
\begin{enumerate}
    \item Train LightGBM ensemble (50 seeds) on training data
    \item Calibrate conformal predictor on 50\% of validation set
    \item Evaluate coverage on held-out validation and test sets
\end{enumerate}

\subsection{Ensemble for Variance Quantification}

We train 50 independent models with random seeds 42--91. For each seed $s$: train LightGBM classifier $M_s$, split validation 50/50 into calibration ($\mathcal{D}_{\text{cal}}^s$) and evaluation ($\mathcal{D}_{\text{eval}}^s$), calibrate conformal predictor $\text{CP}_s$ on $\mathcal{D}_{\text{cal}}^s$, and evaluate on both $\mathcal{D}_{\text{eval}}^s$ and the test set. We report mean $\pm$ standard deviation across 50 trials. This is \textit{not} ensemble prediction---each seed is an independent experimental trial.

\subsection{Feature Overlap Metric}

We define feature temporal stability using Jaccard similarity (Equation~\ref{eq:jaccard}):
\begin{equation}\label{eq:jaccard}
J(f) = \frac{|A_{\text{train}} \cap A_{\text{test}}|}{|A_{\text{train}} \cup A_{\text{test}}|}
\end{equation}
where $A_{\text{train}}$ and $A_{\text{test}}$ are the sets of unique values for feature $f$ in train and test data.

\section{Results}\label{sec:results}

\subsection{Coverage Degradation Across Tasks}

Table~\ref{tab:main_results} shows coverage degradation across 8 supply chain regression tasks experiencing severe feature turnover (Jaccard $\approx$0). Coverage drops range from 0\% to 86.7\%---two orders of magnitude variation under identical temporal shift. Table~\ref{tab:overlap} summarizes feature overlap for representative tasks.

\begin{table*}[t]
\centering
\caption{Coverage Degradation Under COVID-19 Distribution Shift (50 Model Seeds).
High-variance tasks ($^*$) show severely skewed distributions (CV $>$ 50\%);
median and IQR provide robust statistics.}\label{tab:main_results}
\small
\begin{tabular}{@{}lcccccc@{}}
\toprule
\multirow{2}{*}{Task} & \multirow{2}{*}{Classes} &
\multicolumn{2}{c}{Val Coverage (\%)} &
\multicolumn{2}{c}{Test Coverage (\%)} &
\multirow{2}{*}{Cat} \\
\cmidrule(lr){3-4} \cmidrule(lr){5-6}
& & Mean$\pm$SD & Med(IQR) & Mean$\pm$SD & Med(IQR) & \\
\midrule
s-shipcond   & 45  & 93.5$\pm$1.0  & 93.5(1.4)  & 21.8$\pm$27.0 & 13.5(39.7) & SEV   \\
s-group      & 459 & 83.6$\pm$6.7  & 86.8(6.8)  & 12.4$\pm$32.3 & \textbf{0.5(0.7)}   & SEV$^*$ \\
s-payterms   & 137 & 90.8$\pm$0.6  & 90.6(0.7)  & 13.7$\pm$27.0 & \textbf{0.1(0.3)}   & SEV   \\
i-plant      & 35  & 92.0$\pm$0.9  & 91.7(1.3)  & 81.4$\pm$8.4  & 81.8(12.7) & ROB   \\
i-shippoint  & 69  & 91.2$\pm$0.8  & 91.1(0.9)  & 72.7$\pm$36.1 & 89.9(7.0)  & SEV$^*$ \\
s-incoterms  & 13  & 95.5$\pm$0.6  & 95.6(0.9)  & 87.0$\pm$8.0  & 92.3(17.3) & ROB   \\
i-incoterms  & 13  & 95.0$\pm$0.4  & 95.0(0.6)  & 83.7$\pm$9.9  & 82.1(19.8) & ROB   \\
s-office     & 25  & 99.9$\pm$0.0  & 99.9(0.0)  & 99.9$\pm$0.0  & 99.9(0.0)  & ROB   \\
\bottomrule
\end{tabular}
\vspace{1mm}

\raggedright
\footnotesize
SEV = Severe ($>$70\% drop), ROB = Robust ($<$15\% drop).
$^*$High model variance (std $>$ 30\%, CV $>$ 50\%) indicates knife-edge regime.
\end{table*}

\subsection{Diagnostic Analysis}

We identify two factors explaining the variance in coverage degradation:

\textbf{Factor 1: Task Complexity (Entropy).} Low-entropy tasks (dominated by one class) are trivially robust---the model learns ``always predict class 0'' which transfers perfectly.

\textbf{Factor 2: Feature Temporal Stability (Jaccard).} Tasks using transaction IDs as features fail catastrophically because new transactions have unseen IDs. Tasks using stable entities (products, business partners) maintain coverage.

\begin{table}[t]
\centering
\caption{Feature Overlap for Primary Features.}\label{tab:overlap}
\small
\begin{tabular}{lccc}
\toprule
Task & Primary Feature & Jaccard & Drop \\
\midrule
s-shipcond & SALESDOC (ID) & 0.02 & 93.1\% \\
s-group & SALESDOC (ID) & 0.02 & 86.7\% \\
i-incoterms & PRODUCT, PARTY & 0.58 & 0.5\% \\
s-office & SALESORG & 0.61 & 0.1\% \\
\bottomrule
\end{tabular}
\end{table}

\subsection{Feature Importance Analysis}\label{sec:shap}

To understand the mechanism underlying catastrophic failures, we conducted SHAP~\citep{lundberg2017unified} analysis on two contrasting tasks: sales-shipcond (catastrophic, 71.6\% drop) and sales-office (robust, 0.1\% drop).

\textbf{Surprising finding:} Both tasks exhibited near-complete feature distribution shift ($\sim$0\% Jaccard across top-5 features), yet their coverage drops differed by 700$\times$. This rules out feature value overlap as the primary mechanism and points to feature importance dynamics.

Figure~\ref{fig:shap} reveals the key distinction. In the catastrophic task (Panel A), the dominant feature experiences a concentrated importance explosion (4.5$\times$ increase) while maintaining its top ranking. In the robust task (Panel B), importance redistributes dramatically across features with complete ranking reshuffle (10--20$\times$ increases distributed across all top-5 features).

\textbf{Main finding (severe shift)}: Among tasks with complete feature turnover ($n=8$, Jaccard $< 0.05$), SHAP concentration predicts vulnerability with Spearman $\rho=0.714$, $p=0.047$ (Table~\ref{tab:stratified_correlation}). Tasks with concentration $>$45\% show mean coverage drop of 78.5\%, while tasks with $\leq$45\% show only 7.6\% (70.9 pp separation).

\textbf{Exploratory validation (moderate shift)}: Analysis of 4 additional tasks from clinical trials and motorsports with moderate feature stability (Jaccard 0.13--0.86) shows \textit{no concentration effect} ($\rho=0.632$, $p=0.368$, n.s.). Feature stability itself determines robustness (Figure~\ref{fig:n12_correlation}).

\textbf{Mechanistic interpretation}: Concentration matters \textit{conditionally}: in severe-shift scenarios (complete feature turnover), single-feature dependence is catastrophic. When features remain stable, models leverage feature values directly, rendering concentration less critical.

\begin{table}[t]
\centering
\caption{Stratified Correlation Analysis by Shift Severity.}\label{tab:stratified_correlation}
\begin{tabular}{@{}lcccc@{}}
\toprule
Group & $n$ & Spearman $\rho$ & $p$-value & Sig. \\
\midrule
Severe shift    & 8 & 0.714 & 0.047 & Yes \\
Moderate shift  & 4 & 0.632 & 0.368 & No \\
\midrule
Combined & 12 & 0.676 & 0.016 & Yes$^\dagger$ \\
\bottomrule
\end{tabular}
\vspace{1mm}

\raggedright
\footnotesize
$^\dagger$Combined correlation conflates two different mechanisms and is misleading.
\end{table}

\begin{figure*}[t]
\centering
\includegraphics[width=0.85\textwidth]{figures/figure3_feature_importance.pdf}
\caption{\textbf{Feature Importance Analysis.} (A) Catastrophic task: dominant feature explodes 4.5$\times$ while maintaining top rank. (B) Robust task: importance redistributes across features with ranking reshuffle. (C) Importance dynamics: catastrophic tasks show concentrated increase; robust tasks distribute increases. (D) Ranking stability: robust tasks allow greater rank changes. Despite $\sim$0\% Jaccard similarity, coverage drops differ by 700$\times$.}\label{fig:shap}
\end{figure*}

\begin{figure}[t]
\centering
\includegraphics[width=\linewidth]{results/figure_n12_correlation.pdf}
\caption{\textbf{SHAP Concentration Predicts Coverage Degradation.} Spearman $\rho=0.714$, $p=0.047$ for 8 severe-shift tasks. Lighter points: 4 moderate-shift tasks (different mechanism). The 40\% threshold applies to severe-shift scenarios.}\label{fig:n12_correlation}
\end{figure}

\section{Extended Experiments}\label{sec:extended}

\subsection{Does Adaptive Conformal Help?}

We test ACI~\citep{gibbs2021adaptive} which updates the quantile online (Table~\ref{tab:aci}):

\begin{table}[t]
\centering
\caption{ACI Does Not Help Under Severe Shift.}\label{tab:aci}
\begin{tabular}{lc}
\toprule
Method & Test Coverage \\
\midrule
Standard Conformal & 0.2\% \\
ACI ($\gamma$=0.001) & 0.0\% \\
ACI ($\gamma$=0.01) & 0.0\% \\
ACI ($\gamma$=0.05) & 0.0\% \\
\bottomrule
\end{tabular}
\end{table}

ACI fails because conformity scores become meaningless when feature overlap is $\sim$0\%---this is not a calibration problem but a fundamental feature distribution shift.

\subsection{Retraining Analysis}

We test four retraining frequencies on catastrophic (sales-shipcond) and robust (sales-office) tasks over 11 months (Feb--Dec 2020).

\textbf{Catastrophic task} (Table~\ref{tab:retrain}, Figure~\ref{fig:retrain}): Quarterly retraining achieves the highest mean coverage (41.1\%), significantly outperforming no retraining (22.2\%, Wilcoxon $p=0.04$), an improvement of 18.9 pp. Monthly retraining also outperforms no retraining (32.0\%) but shows higher variance and occasional collapse (min: 0.6\%). Quarterly strikes optimal cost-effectiveness: 3 retrains/year achieves 85\% of maximum improvement while monthly (10 retrains/year) achieves 54\% at 3$\times$ the cost.

\textbf{Robust task}: Coverage remains 99.8--99.9\% regardless of retraining frequency, confirming distributed feature importance eliminates need for retraining.

\begin{table}[t]
\centering
\caption{Retraining Impact on Catastrophic Task (Sales-Shipcond).}\label{tab:retrain}
\begin{tabular}{lccc}
\toprule
Frequency & Mean Cov. & Min Cov. & Std \\
\midrule
No retrain                &  22.2\% & 10.8\% & 12.1\% \\
Bi-annual$^{*}$           &  27.0\% & 11.4\% &  8.2\% \\
\textbf{Quarterly}$^{\dagger}$ & \textbf{41.1\%} & \textbf{20.7\%} & 23.4\% \\
Monthly$^{\ddagger}$      &  32.0\% &  0.6\% & 28.3\% \\
\bottomrule
\end{tabular}
\vspace{1mm}

\raggedright
\footnotesize
$^{*}$vs no retrain: $p=0.04$.
$^{\dagger}$vs no retrain: $p=0.04$; vs monthly: $p=0.24$ (n.s.).
$^{\ddagger}$vs no retrain: $p<0.05$.
\end{table}

\begin{figure*}[t]
\centering
\includegraphics[width=0.85\textwidth]{results/retraining/retrain_coverage_over_time.pdf}
\caption{\textbf{Retraining restores coverage for catastrophic tasks.} Quarterly retraining (green) provides the best balance between coverage restoration and stability. Monthly (blue) shows high variance. No retraining (red) exhibits severe degradation post-COVID.}\label{fig:retrain}
\end{figure*}

\subsection{Cross-Domain Validation}

We validate on clinical trials and motorsports datasets (Table~\ref{tab:cross}):

\begin{table}[t]
\centering
\caption{Cross-Domain Validation.}\label{tab:cross}
\small
\begin{tabular}{lcccc}
\toprule
Task & Val & Test & Drop & Jaccard \\
\midrule
\multicolumn{5}{l}{\textit{Clinical Trials (Classification)}} \\
study-outcome & 100.0\% & 100.0\% & 0.0\% & -- \\
study-adverse & 88.6\% & 25.5\% & 63.1\% & -- \\
site-success & 94.8\% & 42.8\% & 52.0\% & -- \\
\midrule
\multicolumn{5}{l}{\textit{Regression Tasks}} \\
s-adverse & 91.9$\pm$0.4 & 88.5$\pm$0.4 & 3.5 & 0.87 \\
si-success & 99.5$\pm$0.2 & 99.5$\pm$0.2 & 0.0 & 0.95 \\
d-position & 92.6$\pm$0.4 & 82.6$\pm$0.5 & 10.0 & 0.70 \\
\bottomrule
\end{tabular}
\end{table}

The pattern replicates across domains and task types: high feature overlap (Jaccard $> 0.85$) yields minimal degradation, while moderate overlap shows moderate degradation.

\subsection{Placebo Test}

To establish that COVID-19 represents a unique shift (not normal drift), we run a placebo test using pre-COVID data: train on 2018, validate on 2019-H1, test on 2019-H2 (Table~\ref{tab:placebo} in Appendix). COVID causes 10--200$\times$ more degradation than normal temporal drift for vulnerable tasks (median 43$\times$). Robust tasks show similar small drops in both periods.

\section{Practitioner Decision Framework}\label{sec:framework}

Before deploying conformal prediction in non-stationary settings:

\begin{enumerate}
    \item \textbf{Compute SHAP concentration} = (Top feature importance) / (Total importance)
    \item \textbf{If concentration $\leq$ 40\%}: Task is ROBUST (distributed importance)
    \item \textbf{If concentration $>$ 40\%}: Check for protective factors---secondary features with Jaccard $> 0.5$ AND importance $> 15\%$
    \begin{itemize}
        \item If protective factors exist $\rightarrow$ ROBUST
        \item Otherwise $\rightarrow$ VULNERABLE: implement quarterly retraining
    \end{itemize}
    \item \textbf{For robust tasks}: Skip retraining to save cost
    \item \textbf{Monitor coverage drift}: Track empirical coverage as early warning
\end{enumerate}

\section{Discussion}\label{sec:discussion}

\textbf{Why Set Sizes Decrease.} Counter-intuitively, catastrophic tasks show smaller prediction sets at test time (\eg, sales-shipcond: 7.0 $\rightarrow$ 3.0). The model is \emph{confidently wrong}---it produces small sets that miss the true label.

\textbf{High Variance as Warning.} Tasks with coefficient of variation $>$ 50\% across seeds are in knife-edge regimes where small initialization differences lead to qualitatively different representations. This is an early warning that conformal calibration will be unreliable.

\textbf{Limitations.} Our primary findings are validated on 8 supply chain tasks with severe feature turnover. The 40\% concentration threshold is empirically derived from $n=8$ categorical-feature tasks and should be validated before applying to other feature types (continuous, high-dimensional). With $n=8$, our primary correlation ($\rho=0.714$, $p=0.047$) is marginally significant. Retraining experiments span 11 months during acute COVID disruption; longer-term dynamics remain open. We focus on LightGBM; deep learning models may exhibit different failure modes.

\section{Conclusion}\label{sec:conclusion}

Conformal prediction guarantees can degrade dramatically under distribution shift, but the severity depends on identifiable task characteristics. Our COVID-19 natural experiment reveals:
(1) Among 8 tasks with severe feature turnover, coverage drops range from 0\% to 86.7\%;
(2) catastrophic failures correlate with single-feature dependence ($\rho=0.714$, $p=0.047$), not feature instability;
(3) the concentration mechanism applies specifically to severe-shift scenarios---moderate-shift tasks exhibit a different mechanism driven by feature stability;
(4) quarterly retraining restores catastrophic coverage by 19 pp ($p=0.04$) while robust tasks need no retraining;
(5) ACI does not help when features lack temporal stability.

We provide a decision framework: monitor SHAP concentration before deployment, retrain quarterly if vulnerable, skip retraining if robust.

\begin{contributions}
    C.~Lee conceived the idea, designed the experiments, implemented the code,
    conducted all analyses, and wrote the paper.
\end{contributions}

\begin{acknowledgements}
    This research received no specific grant from any funding agency in the public,
    commercial, or not-for-profit sectors.
\end{acknowledgements}

% References
\bibliography{references}

\newpage
\onecolumn

\title{Conformal Prediction Under Distribution Shift: A COVID-19 Natural Experiment\\(Supplementary Material)}
\maketitle

\appendix

\section{Reproducibility Details}\label{app:reproducibility}

\subsection{LightGBM Hyperparameters}

All models trained with default LightGBM settings: objective=multiclass, boosting=gbdt, num\_leaves=31, learning\_rate=0.05, feature\_fraction=0.8, bagging\_fraction=0.8, bagging\_freq=5, num\_boost\_round=500, early\_stopping\_rounds=50, seeds 42--91 (50 seeds). No task-specific tuning was performed.

\subsection{Conformal Prediction}

Method: APS~\citep{romano2020classification}. Target coverage: $(1-\alpha)=0.9$ (90\%). Calibration: random 50\% split of validation set. Quantile: $q = \min(\lceil(n+1)(1-\alpha)\rceil/n, 1)$.

\subsection{SHAP Computation}

Method: TreeExplainer~\citep{lundberg2020local}. Subsample: 10,000 per split. Aggregation: mean absolute SHAP across classes.

\subsection{Computational Resources}

Standard CPU (8 cores, 8GB RAM). Wall-clock: $\sim$3--4 hours for full 8-task, 50-seed suite.

\section{Placebo Test Details}\label{app:placebo}

\begin{table}[!h]
\centering
\caption{Placebo Test: Pre-COVID vs COVID Degradation.}\label{tab:placebo}
\begin{tabular}{lccc}
\toprule
Task & Placebo & COVID & Ratio \\
\midrule
s-shipcond & 0.5\% & 93.1\% & 0.01$\times$ \\
s-group & 2.0\% & 86.7\% & 0.02$\times$ \\
s-payterms & 0.0\% & 33.8\% & $\sim$0$\times$ \\
i-plant & 1.8\% & 29.1\% & 0.06$\times$ \\
i-shippoint & 1.5\% & 18.9\% & 0.08$\times$ \\
s-incoterms & 1.2\% & 3.6\% & 0.33$\times$ \\
i-incoterms & 0.6\% & 0.5\% & 1.27$\times$ \\
s-office & 0.1\% & 0.1\% & 0.80$\times$ \\
\bottomrule
\end{tabular}
\end{table}

\section{High Model Variance Analysis}\label{app:variance}

Tasks marked with $^*$ in Table~\ref{tab:main_results} exhibit extreme model variance (std $>$ 30\%):
\begin{itemize}
    \item \textbf{s-group}: Mean 12.4\%, median \textbf{0.5\%}. Most models fail catastrophically; outliers inflate the mean.
    \item \textbf{i-shippoint}: Mean 72.7\%, median \textbf{89.9\%}. Most models maintain coverage; failures deflate the mean.
\end{itemize}
These tasks are in knife-edge regimes. Practitioners should monitor model variance across seeds; CV $>$ 50\% warrants ensemble approaches or architectural changes.

\end{document}
